\section{Placement ablations and the fact-count cliff}
\label{sec:placement}

Where memory lives in the stack is itself informative, and connects our store to
the redundancy shared-KV schemes exploit.

\paragraph{Below the cliff, placement is a flat tie.} On a $2\times2$ grid of read
placement (layer~0 only vs.\ every layer) $\times$ write placement (last block vs.\
every layer), an easy 8-fact task is \recall{100} in every cell; a scrambled
control is at chance. Reading or writing at every layer buys nothing over reading
early and writing late, because the layer-0 token key does the matching wherever
the value is read or written, and a separate store per layer is no better than one
shared store read at every layer.

\paragraph{At the cliff, placement is decisive.} The picture changes at the
fact-count breaking point. A single read/write site (the objective placement)
reproduces the cliff---at $128$ symbols, recall falls from \recall{99.7} at 8 facts
to $9.1\%$ at 16 and $4.2\%$ at 32. Reading \emph{and} writing at every layer into
one shared store escapes it \emph{cold}, with no curriculum: \recall{99.4} and
\recall{96.2} at 16 and 32 facts, versus $9.1\%$ and $4.2\%$ for the single site
(\tabref{tab:cliff}). A unique store per layer also escapes but no better and
more expensively (six stores, seed-noisy), so a single shared multi-layer store is
the better design---the same conclusion CSA2 reaches, from the opposite direction:
the layers are redundant, so one shared KV serves them all.

\begin{table}[t]
\centering
\begin{tabular}{lccc}
\toprule
fact count (128 symbols) & single-layer & shared, read+write @ all & per-layer stores \\
\midrule
16 & $9.1\%$ & \recall{99.4} & \recall{92.7} \\
32 & $4.2\%$ & \recall{96.2} & \recall{100} \\
\bottomrule
\end{tabular}
\caption{The fact-count cliff is a property of a \emph{single} read/write site.
Reading and writing at every layer escapes it cold; a single shared store is as
good as or better than per-layer stores, at lower cost.}
\label{tab:cliff}
\end{table}

\paragraph{On real language the cliff is a delayed transition, not a wall.} On
synthetic keys, multi-layer placement is a complete, curriculum-free escape. On
real-language (BPE) keys it clears the cliff \emph{cold} as well, but only once
training runs past a delayed phase transition. A cold multi-layer model on BPE sits
near-zero recall for thousands of steps, then groks: recall is $0.3\%$ at step
$2{,}000$, $8.8\%$ at $4{,}000$, $77.2\%$ at $6{,}000$ and $98.9\%$ at $8{,}000$
(16 facts), converging to \recall{99.7} by $16{,}000$; 32 facts follows the same
curve one notch later and reaches \recall{98.0}. An earlier, shorter run stopped
inside the pre-transition plateau and read out as a partial lift ($53\%\!\to\!87\%$
at 16, $16\%\!\to\!50\%$ at 32), which we had attributed to a second,
embedding-competition bottleneck; with a sufficient budget that residual gap
closes. The warm-start curriculum reaches the same endpoint---\recall{100} at both
16 and 32 facts---in roughly half the steps, because each warm rung starts with the
read/write loop already established. The curriculum is therefore a training-budget
\emph{optimisation}, not a prerequisite: multi-layer placement removes the cliff,
and the embedding competition survives only as the cost that pushes the transition
late and stalls the lowest-fact cold rung.

\paragraph{Why the exact key is load-bearing.} Every escape and every retrieval
depends on the layer-0 exact-token key. Reading at deep layers gains nothing
because deep queries are contextualised states that no longer match the layer-0
keys; a contextual key---needed for pattern separation---collapses matching to
chance; and dentate-gyrus-style separation is architecturally incompatible with
the exact-match property, because the keys it would separate are byte-identical
token embeddings. The exact key is simultaneously the reason the store works and
the constraint that shapes what can be built on it.
